\chapter{三角函数}

\section{任意角与弧度制}

\begin{definition}[任意角]
角可以看成平面内一条射线绕着端点从一个位置旋转到另一个位置所形成的图形。按逆时针方向旋转形成的角叫做\textbf{正角}，按顺时针方向旋转形成的角叫做\textbf{负角}。
\end{definition}

\begin{definition}[弧度制]
长度等于半径长的弧所对的圆心角叫做 $1$ \textbf{弧度的角}，记作 $1$ rad。
\end{definition}

\begin{theorem}[角度与弧度的换算]
\[ 180^\circ = \pi \text{ rad} \]
\[ 1^\circ = \frac{\pi}{180} \text{ rad} \approx 0.01745 \text{ rad} \]
\[ 1 \text{ rad} = \frac{180^\circ}{\pi} \approx 57.3^\circ \]
\end{theorem}

\begin{note}
记忆口诀：\textbf{一百八十度，等于 $\pi$ 弧度}。
\end{note}

\section{任意角的三角函数}

\begin{definition}[三角函数的定义]
设 $\alpha$ 是一个任意角，它的终边与单位圆交于点 $P(x, y)$，则：
\[ \sin\alpha = y, \quad \cos\alpha = x, \quad \tan\alpha = \frac{y}{x} \ (x \neq 0) \]
\end{definition}

\begin{theorem}[三角函数值的符号]
\begin{center}
\begin{tabular}{c|c}
\toprule
象限 & 正弦、余弦、正切的符号 \\
\midrule
第一象限 & $\sin\alpha > 0$，$\cos\alpha > 0$，$\tan\alpha > 0$ \\
第二象限 & $\sin\alpha > 0$，$\cos\alpha < 0$，$\tan\alpha < 0$ \\
第三象限 & $\sin\alpha < 0$，$\cos\alpha < 0$，$\tan\alpha > 0$ \\
第四象限 & $\sin\alpha < 0$，$\cos\alpha > 0$，$\tan\alpha < 0$ \\
\bottomrule
\end{tabular}
\end{center}
\end{theorem}

\begin{note}
记忆口诀：\textbf{一全正，二正弦，三正切，四余弦}——即第一象限全为正，第二象限正弦为正，第三象限正切为正，第四象限余弦为正。
\end{note}

\section{同角三角函数的基本关系}

\begin{theorem}[平方关系]
\[ \sin^2\alpha + \cos^2\alpha = 1 \]
\end{theorem}

\begin{theorem}[商数关系]
\[ \tan\alpha = \frac{\sin\alpha}{\cos\alpha} \quad (\cos\alpha \neq 0) \]
\end{theorem}

\begin{example}
已知 $\sin\alpha = \dfrac{3}{5}$，且 $\alpha$ 是第二象限角，求 $\cos\alpha$ 和 $\tan\alpha$。
\end{example}

\textbf{解.}

由 $\sin^2\alpha + \cos^2\alpha = 1$，得：
\[ \cos^2\alpha = 1 - \sin^2\alpha = 1 - \frac{9}{25} = \frac{16}{25} \]

因为 $\alpha$ 是第二象限角，所以 $\cos\alpha < 0$。

\[ \cos\alpha = -\frac{4}{5} \]

\[ \tan\alpha = \frac{\sin\alpha}{\cos\alpha} = \frac{3/5}{-4/5} = -\frac{3}{4} \]
\hfill $\square$

\section{诱导公式}

\begin{theorem}[诱导公式]
\begin{enumerate}
    \item 公式一（周期）：$\sin(\alpha + 2k\pi) = \sin\alpha$，$\cos(\alpha + 2k\pi) = \cos\alpha$（$k \in \mathbb{Z}$）
    \item 公式二（负角）：$\sin(-\alpha) = -\sin\alpha$，$\cos(-\alpha) = \cos\alpha$
    \item 公式三（$\pi - \alpha$）：$\sin(\pi - \alpha) = \sin\alpha$，$\cos(\pi - \alpha) = -\cos\alpha$
    \item 公式四（$\pi + \alpha$）：$\sin(\pi + \alpha) = -\sin\alpha$，$\cos(\pi + \alpha) = -\cos\alpha$
    \item 公式五（$\dfrac{\pi}{2} - \alpha$）：$\sin\left(\dfrac{\pi}{2} - \alpha\right) = \cos\alpha$，$\cos\left(\dfrac{\pi}{2} - \alpha\right) = \sin\alpha$
    \item 公式六（$\dfrac{\pi}{2} + \alpha$）：$\sin\left(\dfrac{\pi}{2} + \alpha\right) = \cos\alpha$，$\cos\left(\dfrac{\pi}{2} + \alpha\right) = -\sin\alpha$
\end{enumerate}
\end{theorem}

\begin{note}
记忆口诀：\textbf{奇变偶不变，符号看象限}——
\begin{itemize}
    \item "奇偶"指 $\dfrac{\pi}{2}$ 的奇数倍或偶数倍；
    \item "变"指函数名改变（正弦变余弦，余弦变正弦）；
    \item "符号看象限"指把 $\alpha$ 看作锐角时，原函数在相应象限的符号。
\end{itemize}
\end{note}

\section{三角函数的图象与性质}

\begin{theorem}[正弦函数的图象与性质]
函数 $y = \sin x$：
\begin{itemize}
    \item 定义域：$\mathbb{R}$
    \item 值域：$[-1, 1]$
    \item 周期：$2\pi$
    \item 奇偶性：奇函数
    \item 单调递增区间：$\left[-\dfrac{\pi}{2} + 2k\pi, \dfrac{\pi}{2} + 2k\pi\right]$（$k \in \mathbb{Z}$）
    \item 单调递减区间：$\left[\dfrac{\pi}{2} + 2k\pi, \dfrac{3\pi}{2} + 2k\pi\right]$（$k \in \mathbb{Z}$）
\end{itemize}
\end{theorem}

\begin{theorem}[余弦函数的图象与性质]
函数 $y = \cos x$：
\begin{itemize}
    \item 定义域：$\mathbb{R}$
    \item 值域：$[-1, 1]$
    \item 周期：$2\pi$
    \item 奇偶性：偶函数
    \item 单调递增区间：$[\pi + 2k\pi, 2\pi + 2k\pi]$（$k \in \mathbb{Z}$）
    \item 单调递减区间：$[2k\pi, \pi + 2k\pi]$（$k \in \mathbb{Z}$）
\end{itemize}
\end{theorem}

\section{两角和与差的三角函数}

\begin{theorem}[两角和与差的公式]
\begin{align*}
\sin(\alpha \pm \beta) &= \sin\alpha\cos\beta \pm \cos\alpha\sin\beta \\
\cos(\alpha \pm \beta) &= \cos\alpha\cos\beta \mp \sin\alpha\sin\beta \\
\tan(\alpha \pm \beta) &= \frac{\tan\alpha \pm \tan\beta}{1 \mp \tan\alpha\tan\beta}
\end{align*}
\end{theorem}

\begin{theorem}[二倍角公式]
\begin{align*}
\sin 2\alpha &= 2\sin\alpha\cos\alpha \\
\cos 2\alpha &= \cos^2\alpha - \sin^2\alpha = 2\cos^2\alpha - 1 = 1 - 2\sin^2\alpha \\
\tan 2\alpha &= \frac{2\tan\alpha}{1 - \tan^2\alpha}
\end{align*}
\end{theorem}

\begin{theorem}[降幂公式]
\[ \sin^2\alpha = \frac{1 - \cos 2\alpha}{2}, \quad \cos^2\alpha = \frac{1 + \cos 2\alpha}{2} \]
\end{theorem}

\begin{example}
求 $\sin 75^\circ$ 的值。
\end{example}

\textbf{解.}

\[ \sin 75^\circ = \sin(45^\circ + 30^\circ) = \sin 45^\circ\cos 30^\circ + \cos 45^\circ\sin 30^\circ \]

\[ = \frac{\sqrt{2}}{2} \cdot \frac{\sqrt{3}}{2} + \frac{\sqrt{2}}{2} \cdot \frac{1}{2} = \frac{\sqrt{6} + \sqrt{2}}{4} \]
\hfill $\square$
